\chapter{The Discoverable Recursive Method}

This practice is not a fixed sequence of spreads. It is an open instrument in
which any card or constellation can become a new question.

\section{The parent and the child}

The first draw is the parent reading. A card, pair, row, or whole tableau may
become the parent object of a child reading. The child may use the same deck,
another deck, dice, runes, bones, geomantic figures, an I Ching cast, a book,
or another divinatory instrument whose native grammar the reader understands.

The child does not erase the parent. It adds a new layer. The reader records
what was expanded, why it was expanded, which instrument was used, and whether
the new result clarified, complicated, contradicted, or left the parent open.

\section{Bidirectional recursion}

An English card may be inspected by a Marseille Minor. A Marseille Minor may
be grounded by an English card. A Cicero reading may inspect one position in a
playing-card line, and a playing-card line may provide worldly detail for a
tarot constellation. The direction of inspection is chosen by the question.

This is not a claim that all systems secretly mean the same thing. It is a
method for allowing distinct symbolic grammars to interrogate one another while
preserving their native meanings.

\section{Depth and closure}

Recursion must have a stopping condition. The reader may stop when the question
has become clear, when a practical action is visible, when the instrument asks
for rest, when repetition begins, or when the reading reaches a settled basin.
An answer does not become more true merely because more cards have been drawn.

The practice has a symbolic \texttt{T\_STOP}: the point at which the current
reading is complete enough to be lived rather than further interrogated.

\section{Discovery}

Recursive inspection can make data discoverable that a fixed reading might
leave implicit: a missing domain, a hidden assumption, a practical detail below
an archetypal pattern, an emotional pattern beneath an ordinary event, or a
branch that only appears when a single card is isolated. The method is a way of
asking better questions, not a license to draw without end.
